\chapter{Sectional Curvature Comparison II}
\label{chap:pet-sectional-curvature-comparison-ii}

Lower sectional-curvature bounds yield global comparison, critical-point,
topological, and finiteness results.

\section{Critical Point Theory}

The smooth deformation lemma extends to distance functions after regularity is
expressed in terms of minimizing directions.

\begin{lemma}[deformation retraction along a gradient flow]
  \label{lem:pet-ch12-smooth-deformation-retraction}
  \lean{PetersenLib.gradientFlowDeformationRetraction}
  \source{petersen:page-0457,petersen:page-0458}
  \dcref{ch12:12.1.1}
  Let $(M,g)$ be a Riemannian manifold and $f:M\to\mathbb{R}$ a proper smooth
  function. If $f$ has no critical values in the closed interval $[a,b]$, then
  $f^{-1}((-\infty,b])$ and $f^{-1}((-\infty,a])$ are diffeomorphic; more
  precisely, there is a deformation retraction of $f^{-1}((-\infty,b])$ onto
  $f^{-1}((-\infty,a])$, so the inclusion
  $f^{-1}((-\infty,a])\hookrightarrow f^{-1}((-\infty,b])$ is a homotopy
  equivalence.
\end{lemma}
\begin{proof}
  \uses{lem:pet-ch12-smooth-deformation-retraction}
  \sketch
  Choose a compactly supported cutoff $\psi$ equal to $1$ on
  $f^{-1}([a,b])$, and set
  \[
    X=-\psi\cdot\frac{\nabla f}{|\nabla f|^2}.
  \]
  Its complete flow $F^t$ decreases $f$ at unit speed between levels $a$ and
  $b$. Consequently
  \[
    r_t:f^{-1}((-\infty,b])\to f^{-1}((-\infty,b]),\qquad
    r_t(p)=\begin{cases}p&f(p)\le a\\F^{t(f(p)-a)}(p)&a\le f(p)\le b\end{cases}
  \]
  retracts the upper sublevel onto the lower one, while $F^{b-a}$ gives the
  stated diffeomorphism. Properness supplies the compact support needed for a
  complete flow.
\end{proof}

We generalize \cref{lem:pet-ch12-smooth-deformation-retraction} to distance
functions, which need not be $C^1$. Suppose $(M,g)$ is complete and
$K\subset M$ is compact; then $r(x)=|xK|=\min\{|xp|\mid p\in K\}$ is proper.
Wherever $r$ is smooth it has unit gradient and hence is noncritical there,
but $r$ may also have local maxima, where one still wants a notion of
noncriticality. This is provided by the following generalized gradient.

\begin{definition}[segment direction sets; regular, critical, and
  $\alpha$-regular points]
  \label{def:pet-ch12-alpha-regular}
  \lean{PetersenLib.SegmentDirectionSet,PetersenLib.IsAlphaRegularAt,PetersenLib.IsRegularAt}
  \source{petersen:page-0458}
  Let $(M,g)$ be complete, $K\subset M$ compact, and $r(x)=|xK|$. For
  $x\in M$ let $\overrightarrow{xK}\subset T_xM$ denote the set of initial
  unit velocities of unit-speed minimizing segments from $x$ to $K$ (i.e.\ of
  length $r(x)$); when $K=\{q\}$ this is written $\overrightarrow{xq}$. If $r$
  is smooth at $x$ then $\overrightarrow{xK}$ has a single element and
  $-\nabla r(x)=\overrightarrow{xK}$. In general $x$ is \emph{regular}, or
  \emph{noncritical}, for $r$ if $\overrightarrow{xK}$ lies in an open
  hemisphere of the unit sphere of $T_xM$, equivalently there is a unit
  $v\in T_xM$ with $\angle(v,\overrightarrow{xK})<\pi/2$ for every
  $\overrightarrow{xK}\in\overrightarrow{xK}$ (any such $v$ is the center of
  such a hemisphere); otherwise $x$ is a \emph{critical point} of $r$. More
  generally, for $\alpha\in(0,\pi]$, $x$ is \emph{$\alpha$-regular} for $r$ if
  there is a unit $v\in T_xM$ with $\angle(v,\overrightarrow{xK})<\alpha$ for
  every $\overrightarrow{xK}\in\overrightarrow{xK}$ (so regular means
  $\pi/2$-regular), and $R_\alpha(x,K)\subset T_xM$ is the set of all such
  unit directions $v$ at an $\alpha$-regular point $x$.
\end{definition}

It was Berger who first showed that a local maximum of $r$ must be critical
in this sense; his result is a consequence of the following proposition.

\begin{proposition}[properties of regular points]
  \label{prop:pet-ch12-regular-point-properties}
  \lean{PetersenLib.regularPointProperties}
  \source{petersen:page-0459,petersen:page-0460}
  \uses{def:pet-ch12-alpha-regular}
  \dcref{ch12:12.1.2}
  Suppose $(M,g)$ is complete, $K\subset M$ compact, $r(x)=|xK|$. Then:
  \begin{enumerate}
    \item $\overrightarrow{xK}$ is closed, hence compact, for every $x\in M$.
    \item The set of $\alpha$-regular points is open in $M$.
    \item $R_\alpha(x,K)$ is convex for every $\alpha\le\pi/2$.
    \item If $U\subset M$ is open and every point of $U$ is $\alpha$-regular,
      there is a unit vector field $X$ on $U$ with $X|_x\in R_\alpha(x,K)$ for
      all $x\in U$; moreover, if $c$ is an integral curve of $X$ and $s<t$,
      \[
        |c(s)K|-|c(t)K|\ge\cos(\alpha)(t-s).
      \]
  \end{enumerate}
\end{proposition}
\begin{proof}
  \uses{prop:pet-ch12-regular-point-properties}
  \sketch
  Compactness follows by taking limits of minimizing segments. If
  $x_i\to x$ through non-$\alpha$-regular points, parallel limits of a
  direction and of corresponding minimizing segments show that every unit
  vector at $x$ makes angle at least $\alpha$ with some element of
  $\overrightarrow{xK}$; hence non-regularity is closed. For
  $\alpha\le\pi/2$, $R_\alpha(x,K)$ is an intersection of convex open cones.
  Local choices in these cones combine by a partition of unity and
  normalization to a global unit field on any open regular region.

  Along an integral curve of this field, $r$ is absolutely continuous. Vary a
  minimizing segment from $K$ to the curve; Cauchy--Schwarz gives the support
  inequality
  \[
    \tfrac12\big(r(c(s))\big)^2
    \le\tfrac12\Big(\int_0^1|\partial_t\tilde c|\,dt\Big)^2
    \le E(\tilde c_s),
  \]
  with equality at the base point. The first-variation formula therefore gives
  \[
    -\frac{d}{ds}|c(s)K|=
    \cos\angle\bigl(X,\overrightarrow{xK}\bigr)>\cos\alpha
  \]
  almost everywhere. Integrating yields
  $|c(s)K|-|c(t)K|\ge\cos(\alpha)(t-s)$.
\end{proof}

We can now generalize \cref{lem:pet-ch12-smooth-deformation-retraction} to
distance functions.

\begin{lemma}[Grove--Shiohama retraction lemma]
  \label{lem:pet-ch12-retraction-lemma}
  \lean{PetersenLib.grovShiohamaRetraction}
  \source{petersen:page-0461}
  \uses{def:pet-ch12-alpha-regular, prop:pet-ch12-regular-point-properties}
  \dcref{ch12:12.1.3}
  Let $(M,g)$ be complete and $r(x)=|xK|$ the distance to a compact submanifold
  $K\subset M$. If every point of $r^{-1}([a,b])$ is $\alpha$-regular for
  $\alpha<\pi/2$, then $r^{-1}((-\infty,b])$ deformation retracts onto
  $r^{-1}((-\infty,a])$.
\end{lemma}
\begin{proof}
  \uses{lem:pet-ch12-retraction-lemma,prop:pet-ch12-regular-point-properties}
  \sketch
  By \cref{prop:pet-ch12-regular-point-properties}(4), choose a compactly
  supported field $X$ on $r^{-1}([a,b])$ whose flow satisfies
  \[
    r(p)-r(F^t(p))>t\cos(\alpha),\qquad t\ge0,\ p,F^t(p)\in r^{-1}([a,b]).
  \]
  Each $p\in r^{-1}(b)$ reaches level $a$ at a continuous first time
  $t_p\le(b-a)/\cos\alpha$. Hence
  \[
    r_t:r^{-1}((-\infty,b])\to r^{-1}((-\infty,b]),\qquad
    r_t(p)=\begin{cases}p & r(p)\le a\\ F^{t\cdot t_p}(p) & a\le r(p)\le b\end{cases}
  \]
  is a well-defined homotopy from the identity to a retraction of
  $r^{-1}((-\infty,b])$ onto $r^{-1}((-\infty,a])$.
\end{proof}

\begin{remark}[smooth approximation of the retraction]
  \label{rem:pet-ch12-retraction-remark}
  \lean{PetersenLib.grovShiohamaSmoothApproximation}
  \source{petersen:page-0461}
  \uses{lem:pet-ch12-retraction-lemma}
  \dcref{ch12:12.1.4}
  The construction in \cref{lem:pet-ch12-retraction-lemma} also shows that on a
  region where the
  distance function has no critical points it can be approximated by smooth
  functions without critical points.
\end{remark}

\begin{corollary}[diffeomorphism to the normal bundle]
  \label{cor:pet-ch12-normal-bundle-diffeomorphism}
  \lean{PetersenLib.diffeomorphicToNormalBundle}
  \source{petersen:page-0461,petersen:page-0462}
  \uses{lem:pet-ch12-retraction-lemma, def:pet-ch12-alpha-regular}
  \dcref{ch12:12.1.5}
  Suppose $K$ is a compact submanifold of a complete Riemannian manifold $(M,g)$
  and $|xK|$ is regular everywhere on $M-K$. Then $M$ is diffeomorphic to the
  normal bundle of $K$ in $M$. In particular, if $K=\{p\}$, $M$ is diffeomorphic
  to $\mathbb{R}^n$.
\end{corollary}
\begin{proof}
  \uses{cor:pet-ch12-normal-bundle-diffeomorphism}
  \sketch
  Choose a regular field whose negative flow increases $|xK|$ and agrees near
  $K$ with outward normal geodesics. For each unit normal $v$, let $c_v$ be its
  integral curve and put $F(0_p)=p$, $F(tv)=c_v(t)$. Near the zero section this
  is the normal exponential map. Uniqueness of integral curves makes $F$
  injective, while properness of $|xK|$ and monotonicity of the flow make it
  surjective. The normal exponential map and flow maps are local
  diffeomorphisms, so $F:T^+K\to M$ is a diffeomorphism.
  This realizes \cref{cor:pet-ch12-normal-bundle-diffeomorphism}.
\end{proof}

\section{Distance Comparison}

Hinges and triangles may be degenerate; their sides may be geodesics rather
than minimizing segments when explicitly stated.

\begin{definition}[hinge]
  \label{def:pet-ch12-hinge}
  \lean{PetersenLib.Hinge}
  \source{petersen:page-0462}
  A \emph{hinge} in $(M,g)$ consists of two segments (or geodesics)
  $\overline{px}$ and $\overline{xy}$ that form an interior angle
  $\alpha=\angle(\overrightarrow{xp},\overrightarrow{xy})$ at the common vertex
  $x$, where the specified directions are tangent to the given segments.
\end{definition}

\begin{definition}[triangle]
  \label{def:pet-ch12-triangle}
  \lean{PetersenLib.GeodesicTriangle}
  \source{petersen:page-0462,petersen:page-0463}
  A \emph{triangle} in $(M,g)$ consists of three segments (or geodesics)
  $\overline{xy}$, $\overline{yz}$, $\overline{zx}$ meeting pairwise at the three
  vertices $x,y,z$.
\end{definition}

\begin{lemma}[existence of comparison hinges/triangles]
  \label{lem:pet-ch12-comparison-hinge-triangle-exist}
  \lean{PetersenLib.existsComparisonHingeOrTriangle}
  \source{petersen:page-0463}
  \uses{def:pet-ch12-hinge, def:pet-ch12-triangle}
  \dcref{ch12:12.2.1}
  Suppose $(M,g)$ is complete with $\sec\ge k$. Then for each hinge or triangle
  in $M$ there is a comparison hinge or triangle in the space form $S_k^n$ where
  the corresponding segments have the same lengths and, in the hinge case, the
  angle at the corresponding vertex is the same.
\end{lemma}
\begin{proof}
  \uses{lem:pet-ch12-comparison-hinge-triangle-exist}
  \sketch
  For $k>0$, the diameter bound places every side length in
  $[0,\pi/\sqrt{k}]$. A comparison hinge is obtained by laying off the two side
  lengths in $S_k^n$ at the prescribed angle. For a triangle, lay off one side
  and intersect the distance spheres with radii equal to the other two sides;
  the triangle inequalities ensure an intersection. If a side has maximal
  length $\pi/\sqrt{k}$, the diameter rigidity theorem reduces the assertion to
  the model space itself.
\end{proof}

\begin{proposition}[law of cosines in $S_k^n$]
  \label{prop:pet-ch12-law-of-cosines}
  \lean{PetersenLib.lawOfCosinesSpaceForm}
  \source{petersen:page-0464,petersen:page-0465,petersen:page-0466}
  \dcref{ch12:12.2.3}
  Let a triangle in $S_k^n$ have side lengths $a,b,c$ and let $\alpha$ be the
  angle opposite $a$. Then
  \[
    k=0:\ a^2=b^2+c^2-2bc\cos\alpha;\qquad
    k=-1:\ \cosh a=\cosh b\cosh c-\sinh b\sinh c\cos\alpha;
  \]
  \[
    k=1:\ \cos a=\cos b\cos c+\sin b\sin c\cos\alpha.
  \]
\end{proposition}
\begin{proof}
  \uses{prop:pet-ch12-law-of-cosines}
  \sketch
  Along the side of length $c$, the modified distance $\rho=f_k\circ\sigma$
  from the opposite vertex satisfies $\ddot\rho=1-k\rho$ and has initial data
  determined by $b$ and $-\cos\alpha$. For $k=0$,
  $\rho(t)=\tfrac12b^2-b t\cos\alpha+\tfrac12t^2$. Setting $t=c$, $a=|p\sigma(c)|$
  gives $\tfrac12a^2=\tfrac12b^2-bc\cos\alpha+\tfrac12c^2$.
  For $k=-1$, solving with
  $\rho(0)=\cosh b-1$, $\dot\rho(0)=-\sinh b\cos\alpha$ gives
  $\cosh a-1=\cosh b\cosh c-\sinh b\sinh c\cos\alpha-1$.
  For $k=1$, the data
  $\rho(0)=1-\cos b$, $\dot\rho(0)=-\sin b\cos\alpha$ give the spherical
  identity. These are exactly the three displayed formulas.
\end{proof}

\begin{lemma}[Hessian comparison]
  \label{lem:pet-ch12-hessian-comparison}
  \lean{PetersenLib.hessianComparisonSupportSense}
  \source{petersen:page-0466,petersen:page-0467}
  \dcref{ch12:12.2.4}
  Let $(M,g)$ be complete, $p\in M$, $r(x)=|xp|$. If $\sec M\ge k$, then in the
  support sense everywhere $\mathrm{Hess}_kr\le(1-kt)g$, where $\mathrm{Hess}_k$
  denotes the Hessian of the modified distance function $f_k=f_k\circ r$.
\end{lemma}
\begin{proof}
  \uses{lem:pet-ch12-hessian-comparison}
  \sketch
  Jacobi-field Hessian comparison proves the inequality away from the cut
  locus. At a cut point, move the base point slightly along a minimizing
  segment and use the resulting smooth modified distance plus its displacement
  constant as an upper support; letting the displacement tend to zero gives
  the global support-sense inequality.
\end{proof}

\begin{theorem}[Toponogov comparison theorem, 1959]
  \label{thm:pet-ch12-toponogov}
  \lean{PetersenLib.toponogovComparisonTheorem}
  \source{petersen:page-0463,petersen:page-0464,petersen:page-0467,petersen:page-0468,petersen:page-0469,petersen:page-0470}
  \uses{def:pet-ch12-hinge, def:pet-ch12-triangle, lem:pet-ch12-comparison-hinge-triangle-exist}
  \dcref{ch12:12.2.2}
  Let $(M,g)$ be a complete Riemannian manifold with $\sec\ge k$.
  \begin{itemize}
    \item \emph{Hinge version.} For any hinge with vertices $p,x,y\in M$ forming
      angle $\alpha$ at $x$, and any comparison hinge in $S_k^n$ with vertices
      $p_k,x_k,y_k$ (\cref{lem:pet-ch12-comparison-hinge-triangle-exist}),
      $|py|\le|p_ky_k|$.
    \item \emph{Triangle version.} For any triangle in $M$, the interior angles
      are no smaller than the corresponding interior angles of a comparison
      triangle in $S_k^n$.
  \end{itemize}
\end{theorem}
\begin{proof}
  \uses{thm:pet-ch12-toponogov,prop:pet-ch12-law-of-cosines,lem:pet-ch12-hessian-comparison}
  \sketch
  The triangle statement follows from the hinge statement and the model law of
  cosines in \cref{prop:pet-ch12-law-of-cosines}: with two sides fixed,
  increasing the opposite side increases the
  included angle.

  For the hinge, fix a geodesic in $M$ and a model geodesic with prescribed
  initial distance and radial derivative. By
  \cref{lem:pet-ch12-hessian-comparison}, for
  $k=0$ the distance-function difference has nonnegative second derivative,
  while for $k=-1$ it satisfies $\psi''\ge\psi$. For $k=1$, compare with
  $\zeta''=-(1+\varepsilon)\zeta$ and use a support-function maximum argument
  until $\zeta$ vanishes. Letting $\varepsilon$ and the initial perturbation
  tend to zero proves the hinge inequality.
\end{proof}

\begin{remark}[role of segments in the proof]
  \label{rem:pet-ch12-toponogov-segment-remark}
  \lean{PetersenLib.toponogovSegmentRemark}
  \source{petersen:page-0470}
  \uses{thm:pet-ch12-toponogov}
  \dcref{ch12:12.2.5}
  The proof never uses that the objects being compared in $M$ are segments; it
  only requires that the geodesics used in $S_k^n$ are segments. For $k\le0$
  this is automatic. For $k>0$ it means the comparison geodesic must have length
  $<\pi/\sqrt k$, the condition used in the $k=1$ case of the proof.
\end{remark}

\section{Sphere Theorems}

In this section the curvature lower bound is normalized to $\sec\ge1$. Then
$\operatorname{diam}M\le\pi$, with equality only for $S^n(1)$. An odd-dimensional
$M$ is orientable; an orientable even-dimensional $M$ is simply connected and
satisfies $\operatorname{inj}M\ge\pi/\sqrt{\max\sec}$. The same injectivity
estimate holds when $M$ is simply connected and $\max\sec<4$, in which case
$M$ is homotopy equivalent to a sphere.

\begin{theorem}[quarter-pinched sphere theorem, Rauch--Berger--Klingenberg]
  \label{thm:pet-ch12-quarter-pinched-sphere}
  \lean{PetersenLib.quarterPinchedSphereTheorem}
  \source{petersen:page-0471}
  \uses{thm:pet-ch12-toponogov}
  \dcref{ch12:12.3.1}
  If $M$ is a simply connected closed Riemannian manifold with
  $1\le\sec\le4-\delta$, then $M$ is homeomorphic to a sphere.
\end{theorem}
\begin{proof}
  \uses{thm:pet-ch12-quarter-pinched-sphere,thm:pet-ch12-toponogov,def:pet-ch12-alpha-regular}
  \sketch
  The pinching gives $\mathrm{inj}(M)\ge\pi/\sqrt{4-\delta}>\pi/2$. Choose
  diametral $p,q$. Every $x$ lies in one of the balls
  $B(p,\pi/\sqrt{4-\delta})$ and $B(q,\pi/\sqrt{4-\delta})$, as follows from
  the hinge comparison below; the two-ball decomposition gives the sphere.

  Fix $x\in M$ and consider the triangle $p,x,q$. Suppose $|xq|>\pi/2$; we show
  $|px|<\pi/2$. Since $q$ maximizes distance from $p$, $q$ is not a regular
  point for the distance function to $p$, so for a segment $\overline{xq}$
  there is a segment $\overline{pq}$ with interior angle
  $\alpha=\angle(\overrightarrow{qx},\overrightarrow{qp})\le\pi/2$ at $q$. The
  hinge version of \cref{thm:pet-ch12-toponogov} gives
  \[
    \cos|px|\ge\cos|xq|\cos|pq|+\sin|xq|\sin|pq|\cos\alpha\ge\cos|xq|\cos|pq|.
  \]
  Since $|xq|,|pq|>\pi/2$, the right-hand side is positive, forcing
  $\cos|px|>0$, i.e.\ $|px|<\pi/2$. By symmetry the same argument with $p,q$
  interchanged shows that if $|px|>\pi/2$ then $|xq|<\pi/2$, so every $x$ lies
  in one of the two balls.
\end{proof}

\begin{theorem}[Berger 1962 and Grove--Shiohama diameter sphere theorem, 1977]
  \label{thm:pet-ch12-diameter-sphere}
  \lean{PetersenLib.diameterSphereTheorem}
  \source{petersen:page-0472,petersen:page-0473}
  \uses{thm:pet-ch12-toponogov}
  \dcref{ch12:12.3.2}
  If $(M,g)$ is closed with $\sec\ge1$ and $\mathrm{diam}>\pi/2$, then $M$ is
  homeomorphic to a sphere.
\end{theorem}
\begin{proof}
  \uses{thm:pet-ch12-diameter-sphere,thm:pet-ch12-toponogov,def:pet-ch12-alpha-regular}
  \sketch
  Choose diametral \(p,q\) with \(|pq|>\pi/2\). The index argument shows that a geodesic
  loop at \(p\) of length at most \(\pi\) contradicts the hinge inequality,
  so all loops at \(p\) have length greater than \(\pi\); the index argument
  gives the sphere conclusion.

  The hinge inequality also shows that for every \(x\notin\{p,q\}\), the
  angle between segments to \(p\) and \(q\) exceeds \(\pi/2\); otherwise a
  cosine comparison forces a distance larger than the diameter. Thus \(q\) is the
  unique critical point of \( |\cdot p| \). A vector field increasing this
  distance, equal to the corresponding gradients near \(p\) and \(q\), foliates
  the complement of two small balls by flow lines. The resulting product
  \(S^{n-1}\times[0,1]\), glued to the balls, is homeomorphic to \(S^n\).
\end{proof}

\begin{theorem}[almost quarter-pinched CROSS classification]
  \label{thm:pet-ch12-cross-pinching}
  \lean{PetersenLib.crossPinchingClassification}
  \source{petersen:page-0474}
  \dcref{ch12:12.3.3}
  Let $(M,g)$ be simply connected of dimension $n$. There is $\varepsilon(n)>0$
  such that if $1\le\sec\le4+\varepsilon$, then $M$ is diffeomorphic to a sphere
  or to one of the compact rank-one symmetric spaces (CROSS) $\mathbb{CP}^{n/2}$,
  $\mathbb{HP}^{n/4}$, $\mathbb{OP}^2$.
\end{theorem}

\begin{theorem}[classification at diameter at least $\pi/2$]
  \label{thm:pet-ch12-diameter-classification}
  \lean{PetersenLib.diameterHalfPiClassification}
  \source{petersen:page-0474}
  \dcref{ch12:12.3.4}
  If $(M,g)$ is closed with $\sec\ge1$ and $\mathrm{diam}\ge\pi/2$, then one of
  the following holds:
  \begin{enumerate}
    \item $M$ is homeomorphic to a sphere;
    \item $M$ is isometric to a finite quotient $S^n(1)/\Gamma$ with $\Gamma$
      acting reducibly (an invariant subspace exists);
    \item $M$ is isometric to $\mathbb{CP}^m$, $\mathbb{HP}^m$, or
      $\mathbb{CP}^m/\mathbb{Z}_2$ for $m$ odd;
    \item $M$ is isometric to $\mathbb{OP}^2$.
  \end{enumerate}
\end{theorem}

\section{The Soul Theorem}

\begin{lemma}[Gromov's critical point estimate, 1981]
  \label{lem:pet-ch12-gromov-critical-point-estimate}
  \lean{PetersenLib.gromovCriticalPointEstimate}
  \source{petersen:page-0475}
  \uses{thm:pet-ch12-toponogov, def:pet-ch12-alpha-regular}
  \dcref{ch12:12.4.2}
  If $(M,g)$ is a complete open manifold with $\sec\ge0$, then for every $p\in M$
  the distance function $|xp|$ has no critical points outside some ball
  $B(p,R)$. In particular $M$ has the topology of a compact manifold with
  boundary.
\end{lemma}
\begin{proof}
  \uses{lem:pet-ch12-gromov-critical-point-estimate,thm:pet-ch12-toponogov}
  \sketch
  If  critical for 0\cdot p| a segment from   angle at
  most \(\pi/3\) with \(px\), the flat hinge comparison gives
  \[ |xy|^2\le |py|^2+|xp|^2-|py||xp|. \]
  Criticality at \(x\) supplies a segment to \(y\) making angle at most
  \(\pi/2\), hence \( |py|\le2|xp|\). At most \(6^{n-1}\) directions on
  the unit sphere can be pairwise separated by \(\pi/3\). Therefore a chain
  of critical points whose distances from \(p\) double has length at most
  \(6^{n-1}\), which bounds all critical points in a ball.
\end{proof}

\begin{definition}[totally convex]
  \label{def:pet-ch12-totally-convex}
  \lean{PetersenLib.IsTotallyConvex}
  \source{petersen:page-0476}
  A subset $A\subset M$ of a Riemannian manifold is \emph{totally convex} if
  every geodesic in $M$ joining two points of $A$ lies entirely in $A$.
\end{definition}

\begin{example}[souls of the flat cylinder]
  \label{ex:pet-ch12-totally-convex-cylinder}
  \lean{PetersenLib.flatCylinderSouls}
  \source{petersen:page-0476}
  \uses{def:pet-ch12-totally-convex}
  \dcref{ch12:12.4.3}
  On the flat cylinder $\mathbb{R}\times S^1$, every circle $\{p\}\times S^1$ is
  a geodesic and totally convex; no point of $M$ is totally convex, and every
  such circle is a soul.
\end{example}

\begin{example}[soul of a paraboloid of revolution]
  \label{ex:pet-ch12-totally-convex-paraboloid}
  \lean{PetersenLib.paraboloidSoul}
  \source{petersen:page-0476}
  \uses{def:pet-ch12-totally-convex}
  \dcref{ch12:12.4.4}
  Let $(M,g)$ be a smooth rotationally symmetric metric
  $dr^2+\rho^2(r)d\theta^2$ on $\mathbb{R}^2$ with $\rho'<0$ (a paraboloid of
  revolution). Every geodesic from the origin $r=0$ is a ray to infinity, so the
  origin is a totally convex soul; most other points admit geodesic loops.
\end{example}

\begin{lemma}[superlevel sets of concave functions]
  \label{lem:pet-ch12-concave-superlevel-convex}
  \lean{PetersenLib.concaveSuperlevelTotallyConvex}
  \source{petersen:page-0476}
  \uses{def:pet-ch12-totally-convex}
  \dcref{ch12:12.4.5}
  If $f:(M,g)\to\mathbb{R}$ is concave, meaning $\mathrm{Hess}\,f\le0$ weakly
  everywhere, then every superlevel set $A=\{x\in M\mid f(x)\ge a\}$ is totally
  convex.
\end{lemma}
\begin{proof}
  \uses{lem:pet-ch12-concave-superlevel-convex}
  \sketch
  The restriction of a weakly concave function to any geodesic is concave.
  Its minimum on a compact segment occurs at an endpoint, so a segment whose
  endpoints satisfy \(f\ge a\) remains in the superlevel set.
\end{proof}

\begin{lemma}[properness and concavity of the Busemann infimum]
  \label{lem:pet-ch12-busemann-proper-concave}
  \lean{PetersenLib.busemannInfimumProperConcave}
  \source{petersen:page-0477}
  \uses{def:pet-ch12-totally-convex, lem:pet-ch12-concave-superlevel-convex}
  \dcref{ch12:12.4.6}
  Let $(M,g)$ be complete, noncompact, $\sec\ge0$, and $p\in M$. Let
  $R_p=\{c:[0,\infty)\to M\mid c(0)=p\}$ be the set of rays from $p$, and set
  $f=\inf_{c\in R_p}b_c$, where $b_c$ is the Busemann function of $c$. Then $f$
  is proper and concave.
\end{lemma}
\begin{proof}
  \uses{lem:pet-ch12-busemann-proper-concave,lem:pet-ch12-concave-superlevel-convex}
  \sketch
  Hessian comparison for distance functions in \(\sec\ge0\) shows that each
  Busemann function \(b_c\) is weakly concave. Infima of concave functions are
  concave, hence so is \(f=\inf b_c\).

  If a superlevel set \(A=\{f\ge a\}\) were noncompact, take points of \(A\)
  tending to infinity. Segments from \(p\) to these points converge to a ray
  \(c\). Total convexity of \(A\) places the ray in \(A\), but
  \[ a\le f(c(t))\le b_c(c(t))=-t, \]
  a contradiction. All superlevel sets are compact, so \(f\) is proper.
\end{proof}

\begin{lemma}[interior, boundary, and the supporting hyperplane property]
  \label{lem:pet-ch12-totally-convex-boundary-structure}
  \lean{PetersenLib.totallyConvexBoundaryStructure}
  \source{petersen:page-0477,petersen:page-0478,petersen:page-0479,petersen:page-0480}
  \uses{def:pet-ch12-totally-convex}
  \dcref{ch12:12.4.7}
  If $A\subset(M,g)$ is closed and totally convex, then $A$ has an interior
  $\mathrm{int}\,A$ and a boundary $\partial A=A-\mathrm{int}\,A$, where
  $\mathrm{int}\,A$ is a totally convex submanifold of $M$ with closure $A$.
  Moreover for each $x\in\partial A$ there is an inward vector
  $w\in T_xM$ such that every segment $\overline{xy}$ with $y\in\mathrm{int}\,A$
  satisfies $\angle(w,\overline{xy})<\pi/2$ (the \emph{supporting hyperplane
  property}); in particular the distance function to a subset of
  $\mathrm{int}\,A$ has no critical points on $\partial A$.
\end{lemma}
\begin{proof}
  \uses{lem:pet-ch12-totally-convex-boundary-structure}
  \sketch
  Let \(k\) be the maximal dimension of a submanifold contained in \(A\), and
  let \(N\) be the union of all such \(k\)-submanifolds. Convexity-radius
  neighborhoods show that any point of \(A\) near \(N\) lies in the same
  submanifold; coning a lower-dimensional slice to a nearby point proves that
  \(N\) is dense in \(A\). Thus \(N=\operatorname{int}A\) is an embedded totally
  convex submanifold and \(\partial A=A-N\).

  At a boundary point, the cone of tangent directions entering \(N\) is open in
  the transported tangent \(k\)-plane and contains no antipodal pair, since an
  antipodal pair would put the boundary point in \(N\). At points realizing a
  small nearest-point distance, a convex geodesic ball identifies this cone with
  the half-space \(\{v:\angle(v,-\nabla r_q)<\pi/2\}\). Approximation by such
  points yields the same supporting half-space at every boundary point. The
  supporting direction makes angle less than \(\pi/2\) with every segment to
  the interior, proving the final noncriticality assertion.
\end{proof}

\begin{lemma}[concavity and uniqueness of distance to the boundary]
  \label{lem:pet-ch12-distance-to-boundary-concave}
  \lean{PetersenLib.distanceToBoundaryConcave}
  \source{petersen:page-0480,petersen:page-0481}
  \uses{def:pet-ch12-totally-convex, lem:pet-ch12-totally-convex-boundary-structure}
  \dcref{ch12:12.4.8}
  Let $(M,g)$ have $\sec\ge0$ and let $A\subset M$ be totally convex. The
  function $r:A\to\mathbb{R}$, $r(x)=|x\partial A|$, is concave on $A$. When
  $\sec>0$, any maximum of $r$ is unique.
\end{lemma}
\begin{proof}
  \uses{lem:pet-ch12-distance-to-boundary-concave,lem:pet-ch12-totally-convex-boundary-structure}
  \sketch
  For a segment from a boundary point \(p\) to an interior point \(q\), take
  the hypersurface \(H\) through \(p\) perpendicular to the segment. The
  supporting-hyperplane property makes \( |\cdot H|\) an upper support for
  \(r(x)=|x\partial A|\). Along the segment, the Riccati equation gives
  \[ \frac{d}{dt}\operatorname{Hess}|\cdot H|(E,E)
     =-R(E,\nabla f,\nabla f,E)-\operatorname{Hess}^2f(E,E)\le0, \]
  so the support is concave (strictly when \(\sec>0\)). A hypersurface chosen
  near a possible cut endpoint supplies the same support there. Therefore \(r\)
  is concave. Composing with an increasing concave function preserves
  concavity, and strict concavity in positive curvature gives uniqueness of a
  maximum.
\end{proof}

\begin{theorem}[soul theorem, Gromoll--Meyer 1969 and Cheeger--Gromoll 1972]
  \label{thm:pet-ch12-soul-theorem}
  \lean{PetersenLib.soulTheorem}
  \source{petersen:page-0475,petersen:page-0481,petersen:page-0482}
  \uses{lem:pet-ch12-busemann-proper-concave, lem:pet-ch12-distance-to-boundary-concave, cor:pet-ch12-normal-bundle-diffeomorphism}
  \dcref{ch12:12.4.1}
  If $(M,g)$ is a complete noncompact Riemannian manifold with $\sec\ge0$, then
  $M$ contains a soul $S\subset M$: a closed totally convex submanifold such
  that $M$ is diffeomorphic to the normal bundle of $S$. When $\sec>0$, the soul
  is a point and $M$ is diffeomorphic to $\mathbb{R}^n$.
\end{theorem}
\begin{proof}
  \uses{thm:pet-ch12-soul-theorem,lem:pet-ch12-busemann-proper-concave,lem:pet-ch12-concave-superlevel-convex,lem:pet-ch12-distance-to-boundary-concave,lem:pet-ch12-totally-convex-boundary-structure,cor:pet-ch12-normal-bundle-diffeomorphism}
  \sketch
  Let \(f\) be the proper concave function of
  \cref{lem:pet-ch12-busemann-proper-concave}, and let \(C_1\) be its maximum
  set. It is nonempty and compact; \cref{lem:pet-ch12-concave-superlevel-convex}
  makes it totally convex. If \(C_i\) has boundary, apply
  \cref{lem:pet-ch12-distance-to-boundary-concave} to
  \( |\cdot\partial C_i|\) and define
  \(C_{i+1}\) as its maximum set. The dimensions strictly decrease: equal
  dimensions would make \(C_{i+1}\) open in \(C_i\), contradicting strict
  increase of distance to \(\partial C_i\) along a segment to the boundary.
  After at most \(n\) steps, obtain a totally convex submanifold \(S\) without
  boundary.

  The supporting-hyperplane property in
  \cref{lem:pet-ch12-totally-convex-boundary-structure} shows that
  \( |\cdot S|\) has
  no critical points outside \(S\). The normal-bundle corollary then gives
  \(M\cong\nu S\) by \cref{cor:pet-ch12-normal-bundle-diffeomorphism}. If
  \(\sec>0\), strict concavity makes every maximum set
  a point, so \(S\) is a point and \(M\cong\mathbb{R}^n\).
\end{proof}

\begin{theorem}[Sharafutdinov retraction, submetry, and soul rigidity]
  \label{thm:pet-ch12-soul-uniqueness-submetry}
  \lean{PetersenLib.sharafutdinovSubmetryToSoul}
  \source{petersen:page-0482}
  \uses{thm:pet-ch12-soul-theorem}
  \dcref{ch12:12.4.9}
  Let $S$ be a soul of a complete Riemannian manifold with $\sec\ge0$, arising
  from the construction of \cref{thm:pet-ch12-soul-theorem}.
  \begin{enumerate}
    \item (Sharafutdinov) There is a distance-nonincreasing map
      $\mathrm{Sh}:M\to S$ with $\mathrm{Sh}|_S=\mathrm{id}$; in particular all
      souls of $M$ are isometric to each other.
    \item (Perel'man) $\mathrm{Sh}$ is a submetry; consequently $S$ is a point
      whenever all sectional curvatures based at some single point of $M$ are
      positive.
  \end{enumerate}
\end{theorem}

\section{Finiteness of Betti Numbers}

\begin{lemma}[bounded corank]
  \label{lem:pet-ch12-corank-bound}
  \lean{PetersenLib.corankBound}
  \source{petersen:page-0485,petersen:page-0486,petersen:page-0487}
  \uses{thm:pet-ch12-toponogov, def:pet-ch12-alpha-regular}
  \dcref{ch12:12.5.2}
  Fix a Riemannian $n$-manifold $M$ with $\sec\ge0$. The \emph{corank} of
  $A\subset M$ is the largest $k$ such that there are balls
  $B(p_1,r_1),\dots,B(p_k,r_k)$ with: (a) a critical point $x_i$ for $p_i$ with
  $|p_ix_i|=10r_i$; (b) $r_i\ge3r_{i-1}$ for $i\ge2$; (c)
  $A\subset\bigcap_iB(p_i,r_i)$. The corank of any $A\subset M$ is bounded by
  $100^n$.
\end{lemma}
\begin{proof}
  \uses{lem:pet-ch12-corank-bound,thm:pet-ch12-toponogov}
  \sketch
  Assume a witness sequence of length \(k>100^n\), fix \(z\in A\), and choose
  two critical directions from \(z\) making angle less than \(1/6\). The scale
  conditions give \( |zx_i|\le11r_i\), \( |zx_j|\ge9r_j\), and the hinge
  comparison yields \( |x_ix_j|\le|zx_j|-3|zx_i|/4\). Triangle inequalities
  and criticality at \(x_i\) give the reverse estimate
  \( |zx_j|\le|x_ix_j|+6r_i\), while \( |zx_i|\ge9r_i\) gives
  \( |x_ix_j|+27r_i/4\le|zx_j|\), a contradiction. The packing bound
  \(6^{n-1}<100^n\) proves the corank estimate.
\end{proof}

\begin{lemma}[Mayer--Vietoris rank inequality]
  \label{lem:pet-ch12-mayer-vietoris-rank}
  \lean{PetersenLib.mayerVietorisRankInequality}
  \source{petersen:page-0488,petersen:page-0489,petersen:page-0490}
  \dcref{ch12:12.5.3}
  Let $B_i^j$, $i=1,\dots,m$, $j=0,\dots,n+1$, be bounded open subsets of an
  $n$-manifold $M$ with $\overline{B_i^j}\subset B_i^{j+1}$. Then
  \[
    \mathrm{rank}_*\Big(\bigcup_iB_i^0\subset\bigcup_iB_i^{n+1}\Big)
    \le\sum_{l=0}^{n}\sum_{i_0<\cdots<i_l}
    \mathrm{rank}_*\Big(\bigcap_{s=0}^lB_{i_s}^0\subset\bigcap_{s=0}^lB_{i_s}^{n+1-l}\Big).
  \]
\end{lemma}
\begin{proof}
  \uses{lem:pet-ch12-mayer-vietoris-rank}
  \sketch
  Use the double complex whose \(q\)-column is the singular chain complex
  of \(q+1\)-fold intersections, with vertical boundary \(\partial\) and
  alternating Mayer--Vietoris differential \(\bar\partial\). Its total
  homology agrees with the homology of the union. For a finite-dimensional
  space of total cycles representing the relevant classes, filter by the first
  nonzero column. At filtration step \(l\), the quotient maps into the
  relative homology of \(l+1\)-fold intersections between levels \(l\) and
  \(l+1\). Thus
  \[\dim(Z_k^l/Z_k^{l+1})\le\sum_{i_0<\cdots<i_{k-l}}
  \operatorname{rank}_*\left(\bigcap_{s=0}^{k-l} B_{i_s}^{l}\subset
    \bigcap_{s=0}^{k-l} B_{i_s}^{l+1}\right).\]
  The last filtration term dies in the larger union. Summing over \(l\) and
  degrees \(k\le n\) gives the displayed Mayer--Vietoris rank inequality.
\end{proof}

\begin{lemma}[content ratio bound]
  \label{lem:pet-ch12-content-ratio-bound}
  \lean{PetersenLib.contentRatioBound}
  \source{petersen:page-0491,petersen:page-0492}
  \uses{lem:pet-ch12-corank-bound, lem:pet-ch12-mayer-vietoris-rank}
  \dcref{ch12:12.5.4}
  Let $M$ be a Riemannian $n$-manifold with $\sec\ge0$, let
  $\mathrm{cont}\,B(p,r)=\mathrm{rank}_*(B(p,r/2)\subset B(p,r))$ be the
  \emph{content} of a ball, and let $\mathcal{C}(k)$ be the largest content of
  any ball of corank $\ge k$ (\cref{lem:pet-ch12-corank-bound}). There is
  $C(n)$, depending only on $n$, with $\mathcal{C}(k)\le
  C(n)\,\mathcal{C}(k+1)$.
\end{lemma}
\begin{proof}
  \uses{lem:pet-ch12-content-ratio-bound,lem:pet-ch12-corank-bound,lem:pet-ch12-mayer-vietoris-rank}
  \sketch
  Choose an incompressible ball \(B(p,R)\) realizing \(\mathcal C(k)\). Every
  \(x\in B(p,R/4)\) has a critical point for \(x\) in
  \(B(x,R/2)-B(x,R/10)\); otherwise the inclusion of concentric smaller
  balls would preserve content and contradict incompressibility. Hence all
  sufficiently small balls centered in \(B(p,R/4)\) have corank at least
  \(k+1\). Cover \(B(p,R/5)\) by a uniformly bounded disjointed-ball
  family. Apply \cref{lem:pet-ch12-mayer-vietoris-rank} to the expanded family
  \(B_i^j=B(p_i,10^{j-n-3}R)\). Each nonempty intersection is controlled by a
  ball of corank \(k+1\), so each summand is at most \(\mathcal C(k+1)\).
  The number of summands depends only on \(n\), giving
  \(\mathcal C(k)\le C(n)\mathcal C(k+1)\).
\end{proof}

\begin{theorem}[Gromov's finiteness theorem, 1978 and 1981]
  \label{thm:pet-ch12-betti-number-finiteness}
  \lean{PetersenLib.gromovFinitenessTheorem}
  \source{petersen:page-0483,petersen:page-0484,petersen:page-0492}
  \uses{thm:pet-ch12-toponogov, lem:pet-ch12-corank-bound, lem:pet-ch12-content-ratio-bound}
  \dcref{ch12:12.5.1}
  There is $C(n)$ such that any complete $M$ with $\sec\ge0$ satisfies:
  \begin{enumerate}
    \item $\pi_1(M)$ can be generated by $\le C(n)$ elements;
    \item for every field $\mathbb{F}$,
      $\sum_{i=0}^n b_i(M,\mathbb{F})=\sum_{i=0}^n\dim H_i(M,\mathbb{F})\le C(n)$.
  \end{enumerate}
\end{theorem}
\begin{proof}
  \uses{thm:pet-ch12-betti-number-finiteness,thm:pet-ch12-toponogov,lem:pet-ch12-corank-bound,lem:pet-ch12-content-ratio-bound,lem:pet-ch12-gromov-critical-point-estimate}
  \sketch
  For the fundamental group, choose deck transformations successively with
  shortest displacement from a lift \(p\). Toponogov shows that the resulting
  displacement directions are pairwise separated by at least \(\pi/3\), so
  spherical packing bounds their number by \(6^{n-1}\).

  Let \(K=100^n\) be the corank bound. Iterating
  \cref{lem:pet-ch12-content-ratio-bound}
  gives \(\mathcal C(0)\le C(n)^K\mathcal C(K)\), and maximal-corank balls
  compress to content \(1\). A sufficiently large ball contains all homology
  of \(M\), has corank zero by
  \cref{lem:pet-ch12-gromov-critical-point-estimate}, and its content is
  \(\sum_i b_i(M,\mathbb F)\). Thus both the generator number and total Betti
  number are bounded solely in terms of \(n\).
\end{proof}

\section{Homotopy Finiteness}

\begin{lemma}[existence of a uniform regularity angle]
  \label{lem:pet-ch12-alpha-regular-existence}
  \lean{PetersenLib.uniformAlphaRegularityRadius}
  \source{petersen:page-0493,petersen:page-0494,petersen:page-0495,petersen:page-0496}
  \uses{thm:pet-ch12-toponogov, def:pet-ch12-alpha-regular}
  \dcref{ch12:12.6.2}
  For an integer $n>1$ and $v,D,k>0$, there are
  $\alpha=\alpha(n,D,v,k)\in(0,\pi/2)$ and $\delta=\delta(n,D,v,k)>0$ such that
  if $(M,g)$ is a Riemannian $n$-manifold with $\mathrm{diam}\le D$,
  $\mathrm{vol}\ge v$, $\sec\ge-k^2$, and $p,q\in M$ satisfy $|pq|\le\delta$,
  then $p$ is $\alpha$-regular for $q$ or $q$ is $\alpha$-regular for $p$.
\end{lemma}
\begin{proof}
  \uses{lem:pet-ch12-alpha-regular-existence,thm:pet-ch12-toponogov}
  \sketch
  Assume mutually non-\(\alpha\)-regular points \(p,q\) with \( |pq|\le\delta\).
  Their minimizing directions are \(\pi-\alpha\)-dense in the unit spheres.
  Spherical comparison bounds the two cones of directions avoiding these sets
  by \(v/6\) after choosing \(\alpha\) small; choose \(r\) with
  \(v(n,-1,r)=v/6\). If a point lies outside both cones and both distances
  exceed \(r\), the two hyperbolic hinge inequalities imply
  \[\cosh|xq|\le\cosh|xp|+f(|pq|),\qquad
    \cosh|xp|\le\cosh|xq|+f(|pq|),\]
  where \(f(t)=(\cosh t-1)\cosh D-\sinh t\sinh r\cos\alpha\) has
  \(f(0)=0\), \(f^{\prime}(0)<0\). For small \(\delta\), these inequalities are
  contradictory. Four sets therefore cover \(M\), but their volumes total less
  than \(v\), proving the lemma.
\end{proof}

\begin{corollary}[$\alpha$-regularity of the diagonal]
  \label{cor:pet-ch12-diagonal-alpha-regular}
  \lean{PetersenLib.diagonalAlphaRegular}
  \source{petersen:page-0496}
  \uses{lem:pet-ch12-alpha-regular-existence}
  \dcref{ch12:12.6.3}
  With $\alpha,\delta$ as in \cref{lem:pet-ch12-alpha-regular-existence}, in
  $M\times M$ with the product metric let $\Delta=\{(x,x)\mid x\in M\}$. Any
  point within distance $\delta/\sqrt2$ of $\Delta$ is $\alpha$-regular for
  $\Delta$.
\end{corollary}
\begin{proof}
  \uses{cor:pet-ch12-diagonal-alpha-regular,lem:pet-ch12-alpha-regular-existence}
  \sketch
  A minimizing segment in the product from \((p,q)\) to the diagonal has
  terminal velocities \((v,-v)\), and joining its two components gives a
  minimizing segment from \(p\) to \(q\). Hence \(\sqrt2\,|(p,q)\Delta|=|pq|\).
  The uniform regularity lemma for \(p,q\) transfers through this bijection to
  \(\alpha\)-regularity for the diagonal.
\end{proof}

\begin{corollary}[nearby maps are homotopic]
  \label{cor:pet-ch12-close-maps-homotopic}
  \lean{PetersenLib.closeMapsAreHomotopic}
  \source{petersen:page-0496,petersen:page-0497}
  \uses{cor:pet-ch12-diagonal-alpha-regular, lem:pet-ch12-retraction-lemma}
  \dcref{ch12:12.6.4}
  With $\delta$ as in \cref{lem:pet-ch12-alpha-regular-existence}, if
  $f_0,f_1:X\to M$ are continuous with $|f_0(x)f_1(x)|<\delta$ for all $x\in X$,
  then $f_0,f_1$ are homotopic.
\end{corollary}
\begin{proof}
  \uses{cor:pet-ch12-close-maps-homotopic,cor:pet-ch12-diagonal-alpha-regular,lem:pet-ch12-retraction-lemma}
  \sketch
  By \cref{cor:pet-ch12-diagonal-alpha-regular}, the diagonal neighborhood is
  \(\alpha\)-regular, so the retraction lemma
  supplies a continuous deformation to \(\Delta\). Its flow paths induce
  continuous curves \(H(p,q,t)\) from \(p\) to \(q\) of length at most
  \(C|pq|\), with \(C=\sqrt2/\cos\alpha\). Applying \(H\) pointwise to
  \(f_0(x),f_1(x)\) gives the homotopy.
\end{proof}

\begin{lemma}[extending nearby points to a simplex]
  \label{lem:pet-ch12-simplex-extension}
  \lean{PetersenLib.simplexExtensionLemma}
  \source{petersen:page-0497,petersen:page-0498,petersen:page-0499}
  \uses{cor:pet-ch12-close-maps-homotopic}
  \dcref{ch12:12.6.5}
  With $C=\sqrt2/\cos\alpha$ and $\delta$ as above, suppose
  $p_0,\dots,p_k\in B(p,r)\subset M$ and $2r(C^k-1)/(C-1)<\delta$. Then there is
  a continuous map $f:\Delta^k\to B\big(p,r+2rC(C^k-1)/(C-1)\big)$ with
  $f(e_i)=p_i$, natural with respect to restriction to faces.
\end{lemma}
\begin{proof}
  \uses{lem:pet-ch12-simplex-extension,cor:pet-ch12-close-maps-homotopic}
  \sketch
  Induct on \(k\). Given an extension over \(\Delta^k\), radially project the
  new simplex to its old face and join its image to \(p_{k+1}\) using the
  homotopy \(H\) from \cref{cor:pet-ch12-close-maps-homotopic}:
  \[\widetilde f(x)=H\left(f\left(\frac{x^0}{\sum_{i\le k}x^i},\ldots,
      \frac{x^k}{\sum_{i\le k}x^i}\right),p_{k+1},x^{k+1}\right).\]
  The hypothesis gives the distance bound \(2r(C^{k+1}-1)/(C-1)<\delta\),
  so \(H\) applies. The length estimate yields
  \( |p\widetilde f|\le r+2rC(C^{k+1}-1)/(C-1)\), and the coordinate
  formula is face-compatible.
\end{proof}

\begin{lemma}[homotopy equivalence from Gromov--Hausdorff closeness]
  \label{lem:pet-ch12-gh-close-homotopy-equivalent}
  \lean{PetersenLib.ghCloseImpliesHomotopyEquivalent}
  \source{petersen:page-0499,petersen:page-0500}
  \uses{lem:pet-ch12-simplex-extension, cor:pet-ch12-close-maps-homotopic}
  \dcref{ch12:12.6.6}
  There is $\varepsilon=\varepsilon(n,k,v,D)>0$ such that if $(M,g_1)$,
  $(N,g_2)$ are Riemannian $n$-manifolds with $\mathrm{diam}\le D$,
  $\mathrm{vol}\ge v$, $\sec\ge-k^2$, and Gromov--Hausdorff distance
  $d_{GH}(M,N)<\varepsilon$, then $M$ and $N$ are homotopy equivalent.
\end{lemma}
\begin{proof}
  \uses{lem:pet-ch12-gh-close-homotopy-equivalent,lem:pet-ch12-simplex-extension,cor:pet-ch12-close-maps-homotopic}
  \sketch
  Realize \(M\) and \(N\) as \(\varepsilon\)-close subsets of a common metric
  space and triangulate them with simplices in \(\varepsilon\)-balls. Map
  corresponding vertices within \(\varepsilon\). By
  \cref{lem:pet-ch12-simplex-extension}, the extension applies when
  \(8\varepsilon(C^n-1)/(C-1)<\delta\), producing compatible maps
  \(f:M\to N\) and \(g:N\to M\). The construction gives
  \[|xf(x)|,|yg(y)|\le7\varepsilon+8\varepsilon C\frac{C^n-1}{C-1}.\]
  Choose \(\varepsilon\) further so that
  \(14\varepsilon+16\varepsilon C(C^n-1)/(C-1)<\delta\). Then the nearby-maps
  corollary \cref{cor:pet-ch12-close-maps-homotopic} makes \(f g\) and
  \(g f\) homotopic to the respective identities.
\end{proof}

\begin{theorem}[Grove--Petersen homotopy finiteness theorem, 1988]
  \label{thm:pet-ch12-homotopy-finiteness}
  \lean{PetersenLib.groveePetersenHomotopyFiniteness}
  \source{petersen:page-0493,petersen:page-0501}
  \uses{lem:pet-ch12-gh-close-homotopy-equivalent}
  For an integer $n>1$ and $v,D,k>0$, the class of Riemannian $n$-manifolds with
  $\mathrm{diam}\le D$, $\mathrm{vol}\ge v$, $\sec\ge-k^2$ contains only finitely
  many homotopy types.
\end{theorem}
\begin{proof}
  \uses{thm:pet-ch12-homotopy-finiteness,lem:pet-ch12-gh-close-homotopy-equivalent}
  \sketch
  The diameter and lower sectional-curvature bounds give uniform covering
  numbers, so the class is Gromov--Hausdorff precompact. Cover it by finitely
  many balls of radius \(\varepsilon(n,k,v,D)\) from
  \cref{lem:pet-ch12-gh-close-homotopy-equivalent}. Any two manifolds in one
  ball are within the required threshold and are
  homotopy equivalent; therefore only finitely many homotopy types occur.
\end{proof}

\section{Exercises}

\begin{remark}[totally geodesic hypersurface sphere criterion]
  \label{rem:pet-ch12-ex-12-8-1}
  \lean{PetersenLib.exercise_12_8_1}
  \source{petersen:page-0501}
  \uses{thm:pet-ch12-quarter-pinched-sphere, def:pet-ch12-alpha-regular}
  \dcref{ch12:12.8.1}
  Let $(M,g)$ be a closed simply connected positively curved manifold. Show
  that if $M$ contains a totally geodesic closed hypersurface, then $M$ is
  homeomorphic to a sphere. (Hint: first show the hypersurface is orientable,
  then that the signed distance function to it has only two critical points, a
  maximum and a minimum; this also shows it suffices to assume
  $H^1(M,\mathbb{Z}_2)=0$.)
\end{remark}

\begin{remark}[curvature and splitting around a codimension-one soul]
  \label{rem:pet-ch12-ex-12-8-2}
  \lean{PetersenLib.exercise_12_8_2}
  \source{petersen:page-0502}
  \uses{thm:pet-ch12-soul-theorem}
  \dcref{ch12:12.8.2}
  Let $(M,g)$ be complete noncompact with $\sec\ge0$ and soul $S\subset M$.
  \begin{enumerate}
    \item If $X\in TS$ and $V\in T^\perp S$, show $\sec(X,V)=0$.
    \item If $S$ has codimension $1$ with trivial normal bundle, show
      $(M,g)=(S\times\mathbb{R},g_S+dr^2)$.
    \item If $S$ has codimension $1$ with nontrivial normal bundle, show a
      double cover splits as in (2).
  \end{enumerate}
\end{remark}

\begin{remark}[solution of the comparison ODE]
  \label{rem:pet-ch12-ex-12-8-3}
  \lean{PetersenLib.exercise_12_8_3}
  \source{petersen:page-0502}
  \uses{thm:pet-ch12-toponogov}
  \dcref{ch12:12.8.3}
  Show that the solution of $\ddot\xi=-(1+\varepsilon)\xi$,
  $\xi(0)=\psi(0)>0$, $\dot\xi(0)=\dot\psi(0)>0$ is
  \[
    \xi(t)=\sqrt{\psi(0)^2+\frac{\dot\psi(0)^2}{1+\varepsilon}}\,
    \sin\Big(\sqrt{1+\varepsilon}\,t+\arctan\Big(\frac{\psi(0)\sqrt{1+\varepsilon}}{\dot\psi(0)}\Big)\Big).
  \]
\end{remark}

\begin{remark}[converse to Toponogov comparison]
  \label{rem:pet-ch12-ex-12-8-4}
  \lean{PetersenLib.exercise_12_8_4}
  \source{petersen:page-0502}
  \uses{thm:pet-ch12-toponogov}
  \dcref{ch12:12.8.4}
  Show that the converse of Toponogov's theorem holds: if for some $k$ the
  conclusion of \cref{thm:pet-ch12-toponogov} holds when hinges (or triangles)
  are compared to the same objects in $S_k^2$, then $\sec\ge k$.
\end{remark}

\begin{remark}[local Toponogov comparison]
  \label{rem:pet-ch12-ex-12-8-5}
  \lean{PetersenLib.exercise_12_8_5}
  \source{petersen:page-0502}
  \uses{thm:pet-ch12-toponogov}
  \dcref{ch12:12.8.5}
  Let $(M,g)$ be complete. Show that if all sectional curvatures on $B(p,2R)$
  are $\ge k$, then Toponogov's comparison theorem
  (\cref{thm:pet-ch12-toponogov}) holds for hinges and triangles in $B(p,R)$.
\end{remark}

\begin{remark}[gradient bound for a distance difference]
  \label{rem:pet-ch12-ex-12-8-6}
  \lean{PetersenLib.exercise_12_8_6}
  \source{petersen:page-0502}
  \uses{thm:pet-ch12-toponogov}
  \dcref{ch12:12.8.6}
  Let $(M,g)$ be complete with $\sec\ge k$ and $f(x)=|xq|-|xp|$, with $|xq|$,
  $|xp|$ smooth at $x$. Show that \cref{thm:pet-ch12-toponogov} bounds
  $|\nabla f|_x$ from below in terms of the distance from $x$ to a segment from
  $q$ to $p$.
\end{remark}

\begin{remark}[volume of a normal tube]
  \label{rem:pet-ch12-ex-12-8-7}
  \lean{PetersenLib.exercise_12_8_7}
  \source{petersen:page-0502,petersen:page-0503}
  \uses{thm:pet-ch12-toponogov, lem:pet-ch12-alpha-regular-existence}
  \dcref{ch12:12.8.7}
  Let $c\subset(M,g)$ be a geodesic in a Riemannian $n$-manifold with
  $\sec\ge-k^2$, and let $T(c,R)$ be the normal tube of radius $R$ around $c$
  (points joined to $c$ by a perpendicular segment of length $\le R$). In
  adapted tube coordinates $(r,s,\theta)$ the metric satisfies, as $r\to0$,
  \[
    g(r)=\begin{pmatrix}1&&\\&\ddots&\\&&0\end{pmatrix}
    +\begin{pmatrix}0&&\\&\ddots&\\&&1\end{pmatrix}r^2+O(r^3),
  \]
  with the first matrix having a single $0$ in the $s$-direction entry and the
  second a single $0$ there. Use the lower curvature bound to bound the volume
  density on the tube from above, and conclude
  $\mathrm{vol}\,T(c,R)\le f(n,k,R,L(c))$ for a continuous $f$ with
  $f\to0$ as $L(c)\to0$. Use this to prove Lem.\ 11.4.9 of the source and
  \cref{lem:pet-ch12-alpha-regular-existence}, showing Toponogov's theorem is
  not needed for the latter.
\end{remark}

\begin{remark}[nonnegative curvature on bundles over $S^2$]
  \label{rem:pet-ch12-ex-12-8-8}
  \lean{PetersenLib.exercise_12_8_8}
  \source{petersen:page-0503}
  \uses{thm:pet-ch12-soul-theorem}
  \dcref{ch12:12.8.8}
  Show that any vector bundle over $S^2$ admits a complete metric of
  nonnegative sectional curvature. (Hint: $k$-dimensional bundles over $S^2$
  are classified by $\pi_1(SO(k))$; there is one $1$-dimensional bundle, the
  $2$-dimensional bundles are parametrized by $\mathbb{Z}$, and for $k>2$ there
  are two $k$-dimensional bundles.)
\end{remark}

\begin{remark}[Busemann convexity from Toponogov comparison]
  \label{rem:pet-ch12-ex-12-8-9}
  \lean{PetersenLib.exercise_12_8_9}
  \source{petersen:page-0503}
  \uses{thm:pet-ch12-toponogov}
  \dcref{ch12:12.8.9}
  Use Toponogov's theorem (\cref{thm:pet-ch12-toponogov}) to show that the
  Busemann function $b_r$ is convex when $\sec\ge0$.
\end{remark}
